\section{خلاصة الفصل}


تم في هذا الفصل بناء النموذج الديناميكي للذراع الروبوتية ذات خمس درجات حرية، وذلك بهدف تحديد العلاقة بين حركة المفاصل والعزوم المطلوبة لتشغيلها. اعتمدت عملية النمذجة على اعتبار الذراع منظومة ميكانيكية متعددة الوصلات، حيث ترتبط الوصلات فيما بينها من خلال مفاصل دورانية، وتمثل إحداثيات هذه المفاصل المتغيرات العامة للنظام. وبذلك يمكن وصف وضع الذراع وسلوكه الديناميكي اعتمادًا على متجه الإحداثيات العامة:

\[
q =
\begin{bmatrix}
\theta_1 &
\theta_2 &
\theta_3 &
\theta_4 &
\theta_5
\end{bmatrix}^{T}
\]

بدأت النمذجة من التوصيف الهندسي والحركي للذراع باستخدام تحويلات الإحداثيات ومصفوفات التحويل المتجانسة، وذلك من أجل تحديد مواضع الوصلات ومراكز كتلها بالنسبة إلى إطار القاعدة. وتعد هذه الخطوة ضرورية قبل الاشتقاق الديناميكي، لأن حساب العطالة وتأثير الجاذبية يعتمد بشكل مباشر على مواضع مراكز الكتل واتجاهات محاور الحركة.

بعد ذلك تم الاعتماد على صياغة لاغرانج لاشتقاق معادلات الحركة. تقوم هذه الصياغة على الفرق بين طاقة الحركة وطاقة الوضع للنظام:

\[
L = K - P
\]

حيث تمثل \(K\) طاقة الحركة الناتجة عن حركة الوصلات، وتمثل \(P\) طاقة الوضع الناتجة عن تأثير الجاذبية. وبما أن الذراع الروبوتية تتكون من عدة وصلات صلبة، فإن طاقة الحركة الكلية يمكن التعبير عنها بدلالة مصفوفة العطالة \(D(q)\)، بينما تعتمد طاقة الوضع على كتل الوصلات ومواضع مراكز كتلها بالنسبة إلى اتجاه الجاذبية.

من خلال تطبيق معادلات أويلر-لاغرانج على متغيرات المفاصل، يتم الحصول على النموذج الديناميكي العام للذراع الروبوتية، والذي يربط عزوم المفاصل بتسارعاتها وسرعاتها وتأثير الجاذبية. وبالصيغة العامة، يمكن كتابة نموذج ديناميك الروبوت بالشكل:

\[
\tau =
D(q)\ddot{q}
+
C(q,\dot{q})\dot{q}
+
D_g(q)
\]

حيث تمثل \(D(q)\) مصفوفة العطالة، ويمثل الحد \(C(q,\dot{q})\dot{q}\) تأثير قوى كوريوليس والقوى النابذة المركزية، بينما يمثل \(D_g(q)\) متجه عزوم الجاذبية.

في هذا العمل، وبهدف تبسيط النموذج والتركيز على العوامل الأكثر تأثيرًا في حساب عزوم التثبيت والحركة الأساسية، تم إهمال قوى كوريوليس والقوى النابذة المركزية. لذلك أصبح النموذج الديناميكي المستخدم لحساب عزوم المفاصل بالشكل:

\[
\tau =
\left(D(q)+I_{act}\right)\ddot{q}
+
D_g(q)
\]

حيث تم إضافة مصفوفة قصور المشغلات \(I_{act}\) إلى مصفوفة العطالة، وذلك لأن المشغلات نفسها تملك قصورًا ذاتيًا يؤثر على العزم المطلوب عند المفاصل. وبما أن قصور المشغل يؤثر على المفصل الموافق له، فقد تمت إضافته إلى العناصر القطرية لمصفوفة العطالة.

من النقاط المهمة التي تم توضيحها أثناء الاشتقاق أن شعاع الجاذبية يجب أن يكون معبرًا عنه بالنسبة إلى إطار القاعدة. فمراكز كتل الوصلات تم تعريفها في الإطارات المحلية الخاصة بها، إلا أن الحد:

\[
U_{pi}r_{cp}
\]

ينتج عن اشتقاق تحويل الوصلة بالنسبة إلى متغير المفصل، ولذلك يعطي تأثير مركز الكتلة معبرًا عنه في إطار القاعدة. وبناءً على ذلك تم اعتماد شعاع الجاذبية بالشكل:

\[
g_0 =
\begin{bmatrix}
0 & 0 & -981 & 0
\end{bmatrix}^{T}
\ \mathrm{cm/s^2}
\]

وقد أدى هذا الاختيار إلى نتائج متوافقة مع التفسير الفيزيائي للذراع، وخصوصًا في حالة المفصل الأول. فعند الوضع الصفري الساكن، كان عزم الجاذبية حول المفصل الأول معدومًا، لأن محور دوران هذا المفصل موازي لاتجاه الجاذبية، وبالتالي فإن وزن الذراع لا يولد عزم قيادة حول هذا المحور.

بعد بناء النموذج، تم الحصول على معادلات العزوم النهائية للمفاصل الخمسة. وقد تبين أن التعويض الكامل في النموذج يؤدي إلى معادلات طويلة ومعقدة، بسبب تعدد الوصلات وكثرة الحدود الناتجة عن مصفوفة العطالة وتحويلات الإحداثيات ومراكز الكتل. لذلك تم عرض الصيغة العامة لمعادلات العزوم داخل متن الفصل، بينما أُرفقت المعادلات التفصيلية الناتجة عن التعويض في ملف خارجي مستقل، وذلك لتجنب إطالة النص والمحافظة على وضوح العرض.

كما تم تحليل عزوم التثبيت الساكنة عند الوضع الصفري، وذلك من خلال تعويض:

\[
q=0,
\qquad
\ddot{q}=0
\]

وفي هذه الحالة تختفي حدود العطالة، وتبقى العزوم الناتجة عن الجاذبية فقط:

\[
\tau = D_g(q)
\]

وقد أظهرت النتائج أن المفاصل القريبة من القاعدة تتحمل عزومًا أكبر نسبيًا، لأنها تتأثر بوزن عدد أكبر من الوصلات اللاحقة وبأذرع عزم أكبر. أما المفاصل القريبة من نهاية الذراع فقد ظهرت عندها عزوم أصغر نسبيًا، لأنها تتحمل أجزاء أقل من الذراع. كما تم تفسير إشارة العزم واتجاهه بالاعتماد على اتجاهات محاور الدوران المعتمدة في نموذج \LR{DH}.

إلى جانب النموذج الرياضي، تم إنشاء نموذج محاكاة للذراع باستخدام \LR{Simscape} \LR{Multibody} اعتمادًا على التصميم الميكانيكي المعد في \LR{SolidWorks}. تم تصدير نموذج التجميع من \LR{SolidWorks} باستخدام إضافة \LR{Simscape} \LR{Multibody} \LR{Link}، ثم تم استيراد ملف \LR{XML} الناتج داخل بيئة \LR{MATLAB} باستخدام الأمر:

\begin{flushleft}
\LR{smimport("Assem1.xml")}
\end{flushleft}

وقد أدى ذلك إلى توليد نموذج \LR{Simulink} يحتوي على الأجسام الصلبة، والإطارات المرجعية، والمفاصل التي تربط أجزاء الذراع مع بعضها. وبذلك أصبح من الممكن تشغيل النموذج داخل بيئة \LR{Simscape} \LR{Multibody} ومشاهدة الذراع كنظام ميكانيكي ثلاثي الأبعاد قابل للمحاكاة.

تكمن أهمية نموذج \LR{Simscape} في أنه يوفر تمثيلًا عدديًا ومرئيًا للذراع، ويمكن استخدامه لدراسة الحركة العامة للنظام وربطه لاحقًا مع نماذج التشغيل أو خوارزميات التحكم داخل \LR{Simulink}. ومع ذلك، تم توضيح أن دقة هذا النموذج تعتمد على صحة نموذج \LR{SolidWorks} الأصلي، وتعريف المواد والكتل ومراكز الكتل، إضافة إلى تمثيل المفاصل ومحاور الدوران بشكل صحيح. كما أن بعض التأثيرات الواقعية، مثل الاحتكاك والخلوص الميكانيكي وحدود المحركات، لا تظهر في النموذج إلا إذا تمت إضافتها يدويًا.

وبذلك يكون هذا الفصل قد جمع بين الاشتقاق الرياضي للنموذج الديناميكي والتمثيل العددي للذراع داخل بيئة المحاكاة. فقد تم اشتقاق معادلات العزوم اعتمادًا على صياغة لاغرانج، وتحديد مصفوفة العطالة ومركبات الجاذبية وقصور المشغلات، ثم تحليل القيم الساكنة الناتجة عند الوضع الصفري، وأخيرًا إنشاء نموذج محاكاة للذراع باستخدام \LR{Simscape} \LR{Multibody}. وتشكل هذه النتائج أساسًا مهمًا للمراحل اللاحقة من العمل، سواء في حساب العزوم المطلوبة، أو اختبار حركة الذراع، أو ربط النموذج مع خوارزميات التحكم.